\minitoc 


% ===================================================================================================================
% Introduction

En statistiques paramétriques, on cherche à déterminer la valeur d'un paramètre 
$ \theta^*$ inconnu. Nous posons ici une autré demarche que celle des estimateurs vue 
précédement. 

% ===================================================================================================================
% Risque 

\section{Risque}

\begin{definition}[Règle de décision]
    On appelle \emph{règle de décision} $ \delta$ une statistique à valeurs dans un
    espace $(\D, \mathcal{D})$ telle que : 
        $$ \delta : ( \mathcal{X}, \mathcal{A})^{\oplus n} \longrightarrow (\D, \mathcal{D}). $$ 
\end{definition}

\begin{remark}
    Une règle de décision est donc une fonction mesurable et indépendante d'un paramètre $ \theta$ puis 
    c'est une statistique. 
\end{remark}

\begin{proposition}
    Il existe différents types de règles de décisions : 
    \begin{itemize}
        \item \textbf{Estimation} : On choisit un paramètre $ \theta$ par un 
        estimateur $ \delta(X_1, \dots, X_n)$ à partir d'informations contenues 
        dans un échantillon. On a donc $ \D = \Theta$ et on peut définir 
        \emph{l'erreur quadratique moyenne} : 
            $$ R_Q( \delta, \theta) = \mathbb{E}_\theta( \| \delta - \theta\|_p^2) $$
        sous $ \myP_\theta$. Ici, on aura $ \D = \mathcal{P}( \Theta)$. 
        \item \textbf{Localisation} : Ici, on choisit plutôt une région pouvant contenir 
        $ \theta$. On veut donc une notion de \emph{confiance} telle que : 
            $$ \forall \theta \in \Theta, \quad \myP_\theta ( \theta \in \delta(X_1, \dots, X_n)) \geqslant 1 - \alpha. $$
        \item \textbf{Test} : On a ici $ \D = \{0,1\}$ ou $ \D = [0,1]$. On cherche 
        à vérifier si une affirmation sur $ \theta$ est vraie plutôt qu'une autre. 
        Elles seront définies par deux types d'hypothèses : \emph{l'hypothèse nulle} $ \mathcal{H}_0$ et 
        \emph{l'hypothèse alternative} $ \mathcal{H}_1$.
        On souhaite donc que : 
            \begin{itemize}
                \item $ \myP( \delta(X_1, \dots, X_n) = 0) $ soit proche de $1$ si $ \theta \in \Theta_0$. 
                \item $ \myP( \delta(X_1, \dots, X_n) = 1) $ soit proche de $1$ si $ \theta \not \in \Theta_0$. 
            \end{itemize} 
    \end{itemize}
\end{proposition}

% ===================================================================================================================
% Estimateurs 

\section{Estimateurs}

\subsection{Définition}

On cherche ici à estimer une transformation $g( \theta)$ d'un paramètre $ \theta$ 
connu. Pour cela, on dispose de : 
    $$ g : 
        \begin{cases}
            \Theta \subseteq \R^p \longrightarrow \R^q \\ 
            \theta \longmapsto g( \theta)    
        \end{cases}
    $$

\begin{definition}[Estimateur]
    On appelle \emph{estimateur} de $g( \theta)$ toute fonction mesurable : 
        $$ \delta : ( \mathcal{X}, \mathcal{A})^{ \oplus n} \longrightarrow (g(\Theta), \mathcal{D}) $$ 
    qui ne dépend pas de $ \theta$. On note $ \Delta$ l'ensemble des estimateurs de 
    $ \theta$. 
\end{definition}

\begin{remark}
    Si $ \hat{ \theta}_n$ est un "bon" estimateur de $ \theta$, alors $ g( \hat{ \theta}_n)$ 
    n'est pas nécessairement un aussi bon estimateur pour $ g( \theta)$. La composition 
    doit s'étudier au cas par cas. 
\end{remark}

\subsection{Comparaisons}

On souhaite qu'un estimateur se comporte comme une boîte noire pointant toujours 
sur $ \theta$ peu importe sa valeur. 

\begin{definition}[Estimateur consistant]
    Une suite d'estimateurs $ \delta_n \in \Delta^\N$ sera dite \emph{consistante}
    si : 
        $$ 
            \forall \theta \in \Theta, \quad 
            \lim_{n \to \infty} \delta_n(X_1, \dots, X_n) = g( \theta) \quad \myP_\theta \text{- p.s}.
        $$
\end{definition}

\begin{definition}[Estimateur Asymptotiquement Gaussien]
    Un suite d'estimateurs $ \delta_n \in \Delta^\N$ sera dite 
    \emph{asymptotiquement gaussienne} si pour tout $ \theta \in \Theta$, il existe 
    une matrice $ V_\theta$ semi-définie positive telle que sous $P_\theta^\infty$, 
    on ait : 
        $$ \sqrt{n} ( \delta_n(X_1, \dots, X_n) - g(\theta)) 
            \underset{n \to \infty}{\overset{\mathcal{N}}{\longrightarrow}} \mathcal{N}(0_p,\V_\theta). 
        $$
\end{definition}

\subsection{Comparaison par le risque}

Un autre façon de comparer deux estimateurs est d'introduire une notion de 
\emph{perte} qui quantifie le degré d'erreur d'estimation de $ g( \theta)$. 
Pour cela, on définit les fonctions de perte ou \emph{loss} en anglais. 

\begin{definition}[Fonction de perte]
    Une fonction de perte $ \ell$ est une fonction mesurable positive : 
        $$ \ell : \Theta \times g(\Theta) \longrightarrow (\R_+, \mathcal{B}_\R) $$ 
    telle que : 
        $$ \boxed{ \forall ( \theta, d) \in \Theta \times g( \Theta), \quad 
            l( \theta, d) = 0 \iff d = g( \theta)    
        }.
        $$
\end{definition}

Comme dit précédemment, on souhaite que nos estimateurs se comportent 
comme des boîtes noires. Il nous faut cela minimiser la perte moyenne 
d'une décision pour toutes les valeurs de $ g(\theta)$ possibles. 
Cette notion sera donc définie comme le \emph{risque}. 

\begin{definition}[Risque]
    On définit la \emph{fonction de risque} associée à une décision $ \delta \in \Delta$, 
    comme l'application : 
        $$ \theta \longmapsto R( \theta, \delta) = \mathbb{E}_\theta( \ell( \theta, \delta)).$$
    Ainsi, un estimateur $ \delta$ sera dit \emph{préférable} à un autre $ \delta' \in \Delta$, ssi : 
        $$ \forall \theta \in \Theta, \quad R( \theta, \delta') \leqslant R( \theta, \delta).$$
    On notera alors $ \delta \succeq \delta'$. 
\end{definition}

\begin{remark}
    Se pose alors un problème fondamental dans la recherche de l'estimateur optimal. 
    Admettons que vous ayez un ami complotiste qui croit les extraterrestres sont 
    parmi nous. Dans tous les mondes qu'il pourrait exister, la plupart ne sont pas 
    dirigés par une secte martienne. Or, il suffit qu'un seul monde le soit pour 
    que vous ne soyez pas meilleur que votre ami pour une valeur spécifique de $ \theta$. 

    C'est exactement le problème que posent les estimateurs constants. On va donc chercher 
    à les exclure en créant des sous-classes $ \Delta_i \subset \Delta$. 
\end{remark}

\begin{proof}
    Supposons qu'il existe $ \delta^* \in \Delta$ un estimateur parfait tel que : 
        $$ \forall \theta \in \Theta, \quad \forall \delta \in \Delta, \quad 
            R(\theta, \delta^*) \leqslant R( \theta, \delta). 
        $$
    alors $ \delta = \delta_{\theta_0}$ qui répond $ \theta_0$ presque-partout 
    quelque soit l'échantillon proposé (un estimateur constant). 
    D'où : 
        $$ R( \theta, \delta_{\theta_0})
            \begin{cases}
                = 0 & \text{ si } \theta = \theta_0 \\ 
                > 0  & \text{ si } \theta \not = \theta_0.
            \end{cases}
        $$ 
    Donc nécessairement $ \forall \theta \in \Theta, R(\theta, \delta^*) = 0$ pour 
    faire mieux que chacun des estimateurs constants. 
\end{proof}

\begin{definition}[Estimateurs admissibles]
    On définit l'ensemble des estimateurs \emph{admissibles}, noté $ \Delta_0$, comme l'ensemble 
    des estimateurs "meilleurs en moyenne". Plus formellement, 
        $$ \Delta_0 = \{ \delta \in \Delta_0 \mid \not \exists \delta' \in \Delta_0, \delta' \succeq \delta\}. $$
\end{definition}

\begin{proposition}
    Supposons que $ \Delta_0$ est stable par combinaison convexe et que $\ell$ est strictement 
    convexe en la décision $ \delta$. Alors si $dd_0$ est admissible dans $ \Delta_0$ 
    et si $ \delta_1$ admet la même fonction de risque, alors $ \delta_1 = \delta_0$ $ \myP_\theta$-p.p 
    pour tout $ \theta \in \Theta$. 
\end{proposition}

\begin{example}
    Soit $ \Delta_0$ la classe des estimateurs sans biais moyen : 
        $$ \Delta_0 := \{ \delta \in \Delta \mid \mathbb{E}_\theta( \delta(X_1, \dots, X_n)) = g( \theta)\}. $$ 
    Soit $ \delta_0, \delta_1 \in \Delta_0$, alors : 
        $$ \forall \alpha \in [0,1], \quad \alpha \delta_0 + (1 - \alpha) \delta_1 \in \Delta_0 $$
    par linéarité de $ \mathbb{E}_\theta(\cdot)$. Donc $ \Delta_0$ est bien stable et : 
        $$ \delta \longmapsto \ell(\theta, \delta) = \| g( \theta) - \delta \|_2^2 $$
    est convexe. Donc dans la classe des estimateurs sans biais, s'il existe un estimateur 
    de plus petit risque, alors il est unique. 
\end{example}

\subsection{Optimalité}

On cherche donc à déterminer des estimateurs optimaux (i.e minimisant la fonction 
de risque). 

\begin{definition}[Estimateur Optimal]
    On dira que $ \delta_{opt}$ est $\ell$-optimal dans $ \Delta_0$ si : 
        $$ \forall \theta \in \Theta, \quad R( \theta, \delta_{opt}) = \min_{ \delta \in \Delta_0} R( \theta, \delta). $$ 
\end{definition}

\begin{definition}[Estimateur Asymptotiquement Optimal]
    On dit qu'une suite $ \delta_n(X_1, \dots, X_n)$ est asymptotiquement 
    $\ell$-optimale relativement à une suite $ \Delta_n$ de classes d'estimateurs si : 
        $$ \forall \theta \in \Theta, \quad \lim_{n \to \infty} 
        \frac{R(\theta, \delta_n(X_1, \dots, X_n))}{\displaystyle\inf_{ \delta \in \Delta_n} R(\theta, \delta)} = 1. $$
\end{definition}

\subsection{Biais}

